\documentclass[french]{book}
\usepackage{teach}
\usepackage{verbatim}
\usepackage{amsmath,amsfonts,amssymb}
\usepackage{pgfplots}
\usepackage{mwe}
\usepackage{yhmath} 
\usepackage{pstricks,pst-plot,pst-text,pst-tree,pst-eucl}

\newcommand{\arc}[1]{\wideparen{#1}}
\RequirePackage{graphicx}
\RequirePackage{ifpdf}
 \ifpdf
   \DeclareGraphicsRule{*}{mps}{*}{}
 \fi

\newcommand{\sysd}[2]{%
       \left\{
       \begin{aligned}
          #1\\
          #2\\
       \end{aligned}
       \right.%
       }
  \usepackage{fancyhdr}

  \usepackage{amsmath}
  \usepackage{amssymb}
  \usepackage{amsfonts}
  \usepackage{amsthm}
  \usepackage{mathrsfs}

  \usepackage{array}
  \usepackage{multicol}
  \setlength{\multicolsep}{3pt}

  \usepackage{graphicx}
  \graphicspath{{Figures/}}
  \usepackage{pstricks,pst-plot,pst-text,pst-tree,pst-eucl}
  \usepackage[all]{xy}

  \usepackage{fancybox}
  \usepackage{color}

%  \usepackage[frenchb]{babel}
  \usepackage{hyperref}

\usepackage{subfig}
\pgfplotsset{compat=newest}
\usepackage[all]{xy}
%\usepackage{geometry}
%\geometry{top=1cm, bottom=1cm, left=2cm, right=2cm}
%\usepackage[utf8]{inputenc}
\usepackage[french]{babel}
\redefinecolor{thm}{title}{bgcolor}{alizarin}
\redefinecolor{prop}{title}{bgcolor}{meatbrown}
\redefinecolor{defn}{title}{bgcolor}{olivine}
\definecolor{alizarin}{rgb}{0.82, 0.1, 0.26}
\definecolor{meatbrown}{rgb}{0.9, 0.72, 0.23}
\definecolor{olivine}{rgb}{0.6, 0.73, 0.45}
\usepackage{mathrsfs}
%%
\newcommand{\R}{\mathbb {R}}
%\newcommand{\C}{\mathds {C}}
\newcommand{\N}{\mathbb {N}}
\newcommand{\Z}{\mathbb {Z}}
\newcommand{\Q}{\mathbb {Q}}
\newcommand{\D}{\mathbb {D}}
%%
\newcommand{\se}{\geq}
\newcommand{\ie}{\leq}
\newcommand{\tm}{\times}
%\newcommand{\ell}{l}
%\newcommand{\ell'}{l'}
%%
\newcommand{\ds}{\displaystyle}
\newcommand{\limn}{\lim_{n\rightarrow +\infty}}
\newcommand{\limo}[1][x]{\ds\lim_{#1\rightarrow 0}}
\newcommand{\limite}[2][x]{\ds\lim_{#1\rightarrow #2}}
\newcommand{\limpinf}[1][x]{\ds\lim_{#1\rightarrow +\infty}}
\newcommand{\limminf}[1][x]{\ds\lim_{#1\rightarrow -\infty}}
\newcommand{\limiteg}[2][x]{\ds\lim_{#1\xrightarrow{<}#2}}
\newcommand{\limited}[2][x]{\ds\lim_{#1\xrightarrow{>}#2}}
%%
\newcommand{\vect}[1]{\overrightarrow{#1}}
\newcommand{\oij}{(\textit{O},{\overrightarrow{i}},{\overrightarrow{j}})}
%
\newcommand{\cp}[2]{%
        \begin{pmatrix}
        #1\\
        #2
        \end{pmatrix}%
        }
%
\newcommand{\calb}{\mathscr{B}}
\newcommand{\calc}{\mathscr{C}}
\newcommand{\cald}{\mathscr{D}}
\newcommand{\cale}{\mathscr{E}}
\newcommand{\calf}{\mathscr{F}}
\newcommand{\calp}{\mathscr{P}}
\newcommand{\cals}{\mathscr{S}}
\newcommand{\calt}{\mathscr{T}}
%
\newcommand{\suite}[1][u]{\left( #1_{n} \right)_{n\in\N}}
\newcommand{\suitep}[1][u]{\left( #1_{n} \right)_{n\in\N^{*}}}
\usetikzlibrary{arrows}

\newcommand{\poly}[1][a]{#1_0+#1_1x+\cdots+#1_nx^n}


\begin{document}
\tableofcontents

\chapter{Applications du produit scalaire}
\section{\'Equations de droites}
\subsection{D\?efinition}
\begin{defn}
Un vecteur non nul $\vect{n}$ est dit normal � une droite $d$ si la direction de $\vect{n}$ est orthogonale � celle de $d$.
\vspace{4mm}
\end{defn}

\subsection{Propri\'et\'es}
\begin{prop}
Le plan est muni d'un rep\`ere orthonormal $\oij$.
\begin{enumerate}
\item Une droite de vecteur normal $\vect{n}\cp{a}{b}$ a une \'equation de la forme $ax+by+c=0$ avec $c\in\R$.
\item \'Etant donn\'es trois r\'eels $a$, $b$ et $c$ o\`u $a$ et $b$ ne sont pas nuls simultan\'ement, l'ensemble des points dont les coordonn\'ees v\'erifient $ax+by+c=0$ est une droite de vecteur normal $\vect{n}\cp{a}{b}$.
\end{enumerate}
\end{prop}

\begin{demo}
1/ Soit $d$ une droite de vecteur normal $\vect{n}\cp{a}{b}$ et soit $A(x_{0};y_{0})\in d$.

Soit $M(x;y)$, on a $\vect{AM}\cp{x-x_{0}}{y-y_{0}}$

$M(x;y)\in d \Leftrightarrow \vect{AM}.\vect{n}=0 \Leftrightarrow (x-x_{0})a+(y-y_{0})b=0 \Leftrightarrow ax+by+c=0$ avec $c=-ax_{0}-by_{0}$.

2/ Soit $d$ l'ensemble des points $M(x;y)$ tels que $ax+by+c=0$ et soit $A(x_{A},y_{A})\in d$.

$M(x;y)\in d \Leftrightarrow ax+by+c=0=ax_{A}+by_{A}+c$

\hspace*{1cm} $\Leftrightarrow a(x-x_{A})+b(y-y_{A})=0 \Leftrightarrow \vect{n}.\vect{AM}=0 \Leftrightarrow \vect{AM}\bot\vect{n}$

\hspace*{1cm} $\Leftrightarrow M$ appartient � la droite passant par $A$ et de vecteur normal $\vect{n}$.
\end{demo}

%\ifthenelse{\value{version}>0}{\newpage}{}
\begin{exemple}
Dans un rep\`ere orthonormal, on consid\`ere les points $A(3;-1)$ et $B(2;4)$. D\'eterminer une \'equation de la m\'ediatrice $m$ de $[AB]$.
\end{exemple}


La m\'ediatrice de [AB] est la droite perpendiculaire � $(AB)$ passant par le milieu $I$ de $[AB]$.

On a $\vect{AB}\cp{-1}{5}$ donc une \'equation de $m$ est de la forme $-x+5y+c=0$.

De plus, $I\left(\dfrac{5}{2};\dfrac{3}{2}\right)\in m$ donc $-\dfrac{5}{2}+5\times \dfrac{3}{2}+c=0$ donc $c=-5$.

Une \'equation de $m$ est donc $-x+5y-5=0$.


\section{\'Equations de cercles}
\subsection{Forme g\'en\'erale}
\begin{prop}
Soit $\oij$ un rep\`ere orthonormal.

Une \'equation du cercle $\calc$ de centre $A(x_{A};y_{A})$ et de rayon $R$ est \[(x-x_{A})^{2}+(y-y_{A})^{2}=R^{2}\]
\end{prop}

\begin{demo}
$M(x;y)\in \calc \Leftrightarrow AM=R \Leftrightarrow AM^{2}=R^{2} \Leftrightarrow (x-x_{A})^{2}+(y-y_{A})^{2}=R^{2}$.
\end{demo}

\begin{exemple}
Soit $\oij$ un rep\`ere orthonormal.

Quelle est la nature de l'ensemble $\cale$ des points $M(x;y)$ tels que $x^{2}+y^{2}-6x+2y+5=0$ ?
\end{exemple}


\begin{align*}
M(x;y)\in\cale&\Leftrightarrow x^{2}+y^{2}-6x+2y+5=0 \Leftrightarrow x^{2}-6x+y^{2}+2y+5=0\\
&\Leftrightarrow (x-3)^{2}-9+(y+1)^{2}-1+5=0 \Leftrightarrow (x-3)^{2}+(y+1)^{2}=5
\end{align*}
$\cale$ est donc le cercle de centre $C(3;-1)$ et de rayon $\sqrt{5}$.


\subsection{Cercle de diam\`etre donn\'e}
\begin{prop}
On consid\`ere deux points $A$ et $B$ du plan. Le cercle $\calc$ de diam\`etre $[AB]$ est l'ensemble des points $M$ du plan tels que $\vect{MA}.\vect{MB}=0$.
\end{prop}

\begin{demo}
$M\in\calc \Leftrightarrow \left|
% use packages: array
\begin{array}{cl}
 & M=A \\
\text{ou} & M=B \\
\text{ou} & AMB \text{ est un triangle rectangle en } M
\end{array}\right.\Leftrightarrow \vect{MA}.\vect{MB}=0$.

Effectivement, soit I un point du plan alors par la relation de Chasles on a: 
$$\vect{MA}.\vect{MB}=(\vect{MI}+\vect{IA}).(\vect{MI}+\vect{IB})$$
Supposons que I soit le milieu de $[AB]$, par suite $\vect{IB}=\vect{AI}=-\vect{IA}$.
$$\vect{MA}.\vect{MB}=(\vect{MI}+\vect{IA}).(\vect{MI}-\vect{IA})$$
On utilise l'identit\'e "remarquable" $(\vect{u}+\vect{v}).(\vect{u}-\vect{v})= \vect{u}^2-\vect{v}^2$)
\begin{align*}
\vect{MA}.\vect{MB}&=(\vect{MI}+\vect{IA}).(\vect{MI}-\vect{IA})\\
&=\vect{MI}^2-\vect{IA}^2
\end{align*}

On conclut de la mani\`ere suivante, si $ \vect{MA}.\vect{MB}=0 $ alors c'est \'equivalent aussi \`a dire que  $\vect{MI}^2-\vect{IA}^2=0$ ou encore � dire $\vect{MI}^2=\vect{IA}^2$.\\

Puisque $MI>0$ et $IA>0$ cela revient � dire que 
$$MI=IA$$

Autrement dit, $[MI]$ est un rayon du cercle $\calc$ de diam�tre $[AB]$, ainsi $ABM$ est incrit dans le cercle et donc rectangle en $M$
\end{demo}

\begin{exemple}
Soit $\oij$ un rep\`ere orthonormal.

D\'eterminer une \'equation du cercle $\calc$ de diam\`etre $[AB]$ avec $A(2;2)$ et $B(6;-2)$.
\end{exemple}


Soit $M(x;y)$. On a $\vect{MA}\cp{2-x}{2-y}$ et $\vect{MB}\cp{6-x}{-2-y}$.

$M(x;y)\in\calc\Leftrightarrow\vect{MA}.\vect{MB}=0 \Leftrightarrow (2-x)(6-x)+(2-y)(-2-y)=0$

\hspace*{2cm} $\Leftrightarrow 12-8x+x^{2}-4-2y+2y+y^{2}=0$

Une \'equation de $\calc$ est donc $x^{2}+y^{2}-8x+8=0$.




\section{Longueurs et angles dans un triangle}
\subsection{Th\'eor\`eme de la m\'ediane}
\begin{minipage}{9.5cm}
\begin{prop}
On consid\`ere deux points $A$ et $B$ du plan et $I$ le milieu de $[AB]$.
Pour tout point $M$ du plan, on a :
\[MA^{2}+MB^{2}=2MI^{2}+\dfrac{1}{2}AB^{2}\]
\end{prop}
\end{minipage}
\hfill
\begin{minipage}{5cm}
\begin{center}
\includegraphics{CoursApplProdScalFigures.1}
\end{center}
\end{minipage}

\begin{demo}
$MA^{2}+MB^{2}=\vect{MA}^{2}+\vect{MB}^{2}=(\vect{MI}+\vect{IA})^{2}+(\vect{MI}+\vect{IB})^{2}$

\hspace*{1cm} $=\vect{MI}^{2}+2\vect{MI}.\vect{IA}+\vect{IA}^{2}+\vect{MI}^{2}+2\vect{MI}.\vect{IB}+\vect{IB}^{2}$

\hspace*{1cm} $=2MI^{2}+2\vect{MI}.(\vect{IA}+\vect{IB})+IA^{2}+IB^{2}$

Or $I$ est le milieu de $[AB]$ donc $IA=IB=\dfrac{1}{2}AB$ donc $IA^{2}=IB^{2}=\dfrac{1}{4}AB^{2}$

De plus $\vect{IA}+\vect{IB}=\vect{0}$.

Ainsi $MA^{2}+MB^{2}=2MI^{2}+2\times \dfrac{1}{4}AB^{2}=2MI^{2}+\dfrac{1}{2}AB^{2}$.
\end{demo}

\begin{exemple}
$ABC$ est un triangle tel que $AB=6$, $AC=8$ et $BC=12$. Calculer $AI$ o� $I$ est le milieu de $[BC]$.
\end{exemple}


D'apr\`es le th\'eor\`eme de la m\'ediane : $AB^{2}+AC^{2}=2AI^{2}+\dfrac{1}{2}BC^{2}$.

On a donc $2AI^{2}=6^{2}+8^{2}-\dfrac{1}{2}\times 12^{2}=28$.

Ainsi $AI=\sqrt{14}$.

\newpage
\subsection{Formules d'Al Kashi}
\begin{minipage}{9.5cm}
\begin{prop}
On consid\`ere un triangle $ABC$. On pose $a=BC$, $b=AC$, $c=AB$, $\widehat{A}=\widehat{BAC}$, $\widehat{B}=\widehat{ABC}$ et $\widehat{C}=\widehat{ACB}$. On a :
\begin{displaymath}
\begin{array}{l}
a^{2}=b^{2}+c^{2}-2bc\cos(\widehat{A}) \\
b^{2}=a^{2}+c^{2}-2ac\cos(\widehat{B}) \\
c^{2}=a^{2}+b^{2}-2ab\cos(\widehat{C})
\end{array}
\end{displaymath}
\end{prop}
\end{minipage}
\hfill
\begin{minipage}{5cm}
\begin{center}
\includegraphics{CoursApplProdScalFigures.2}
\end{center}
\end{minipage}

\begin{demo}
$BC^{2}=\vect{BC}^{2}=(\vect{BA}+\vect{AC})^{2}=\vect{BA}^{2}+2\vect{BA}.\vect{AC}+\vect{AC}^{2}=BA^{2}+AC^{2}-2\vect{AB}.\vect{AC}$

Ainsi $BC^{2}=AB^{2}+AC^{2}-2\times AB\times AC\times \cos(\widehat{BAC})$ soit $a^{2}=b^{2}+c^{2}-2bc\cos(\widehat{A})$
\end{demo}

\begin{exemple}
$ABC$ est un triangle tel que $AC=9$, $AB=5$ et $\widehat{A}=\dfrac{\pi}{3}$. Calculer $BC$.
\end{exemple}


D'apr\`es la formule d'Al Kashi :

$BC^{2}=AB^{2}+AC^{2}-2\times AB\times AC\times \cos(\widehat{A})=5^{2}+9^{2}-2\times 5\times 9\times \cos\left(\dfrac{\pi}{3}\right) =61$

Ainsi $BC=\sqrt{61}$


\subsection{Aire d'un triangle}
\begin{prop}
On consid\`ere un triangle $ABC$ et on appelle $S$ son aire. Avec les notations ci-dessus, on a :
\[S=\dfrac{1}{2}bc\sin(\widehat{A})=\dfrac{1}{2}ac\sin(\widehat{B})=\dfrac{1}{2}ab\sin(\widehat{C})\]
\end{prop}

\begin{minipage}{9.5cm}
\begin{demo}
Soit $H$ le projet\'e orthogonal de $C$ sur $AB$. On a alors $S=\dfrac{1}{2}AB\times CH$.

Si l'angle $\widehat{A}$ est aigu alors $CH=AC\sin(\widehat{A})$.

Si l'angle $\widehat{A}$ est obtus alors $CH=AC\sin(\pi-\widehat{A})=AC\sin(\widehat{A})$

Dans tous les cas $S=\dfrac{1}{2}AB\times AC\times \sin(\widehat{A})$.
\end{demo}
\end{minipage}
\hfill
\begin{minipage}{5cm}
\includegraphics{CoursApplProdScalFigures.3}
\end{minipage}

\subsection{Formule des sinus}
\begin{prop}
On consid\`ere un triangle $ABC$. Avec les notations ci-dessus, on a :
\[\dfrac{a}{\sin(\widehat{A})}=\dfrac{b}{\sin(\widehat{B})}=\dfrac{c}{\sin(\widehat{C})}\]
\end{prop}

\begin{demo}
On a $S=\dfrac{1}{2}bc\sin(\widehat{A})=\dfrac{1}{2}ac\sin(\widehat{B})=\dfrac{1}{2}ab\sin(\widehat{C})$

donc $\dfrac{2S}{abc}=\dfrac{\sin(\widehat{A})}{a}=\dfrac{\sin(\widehat{B})}{b}=\dfrac{\sin(\widehat{C})}{c}$

ainsi $\dfrac{a}{\sin(\widehat{A})}=\dfrac{b}{\sin(\widehat{B})}=\dfrac{c}{\sin(\widehat{C})}$.
\end{demo}

\begin{exemple}
$ABC$ est un triangle tel que $BC=5$, $\widehat{B}=50^{\circ}$ et $\widehat{C}=75^{\circ}$. Calculer $AB$ et $AC$ et donner les valeurs arrondies au dixi\`eme.
\end{exemple}


$\widehat{A}=180^{\circ}-(\widehat{B}+\widehat{C})=55^{\circ}$

$\dfrac{AB}{\sin(\widehat{C})}=\dfrac{AC}{\sin(\widehat{B})}=\dfrac{BC}{\sin(\widehat{A})} \Leftrightarrow \dfrac{AB}{\sin(75^{\circ})}=\dfrac{AC}{\sin(50^{\circ})}=\dfrac{5}{\sin(55^{\circ})}$

On a donc $AB=\dfrac{5\sin(75^{\circ})}{\sin(55^{\circ})}\simeq 5,9$ et $AC=\dfrac{5\sin(50^{\circ})}{\sin(55^{\circ})}\simeq 4,7$.


\end{document}